% Appendix D — N5 wall: lambda-saturation + omega_log ceiling trade-off.
\section{Appendix D --- The N5 wall: $\lambda$-saturation and the
$\omega_{\log}$ ceiling}
\label{app:D}

The central physics finding of the campaign is the \emph{N5 wall}: across
the binary $H_3X$ family, \emph{dynamical stability}, \emph{strong coupling},
and \emph{high $\omega_{\log}$} cannot be achieved simultaneously. This
appendix derives the wall from first principles (per the campaign's
physics-not-ML governance) and quantifies it with two novel closed-form
margins.

\paragraph{First-principles origin.}
In the McMillan--Hopfield decomposition~\cite{mcmillan1968,hopfield1969},
$\lambda = \eta/(M\langle\omega^2\rangle)$, the mean-square phonon frequency
$\langle\omega^2\rangle$ plays a \emph{double role}:
\begin{equation}
\frac{\partial\lambda}{\partial\langle\omega^2\rangle}
  = -\frac{\eta}{M\langle\omega^2\rangle^2} < 0,
\qquad\text{yet}\qquad
\langle\omega^2\rangle > 0 \;\text{required for stability.}
\end{equation}
To maximise $\lambda$ one wants $\langle\omega^2\rangle \to 0^+$ (soft
lattice); but that is precisely the threshold at which a mode crosses into
$\omega^2 < 0$ (imaginary, lattice collapse). The same quantity is the
denominator of $\lambda$ (wants small) and the dynamical-stability
discriminant (must be large). The wall is structural, not accidental.

\paragraph{The wall, drawn.}
\begin{verbatim}
 lambda
 2.5 -  unstable region  (H3O harmonic lambda=2.479, 16/16 imag artifact)
     |   ############ harmonic strong-lambda wall ############
     |    \  SSCHA renormalize => lambda collapses (S = 0.40)
 2.0 -   H3Br *          \
     |   H3Si *           \    stable region
 1.5 -   H3As-R3m *        \   lambda_max ~ 2 ceiling
     |   H3Cl *             \  (cannot enter above Tc=200 K isocontour)
 1.0 -   H3O-anh *----------\----------- Tc = 200 K isocontour (target)
     |   H3F *              |
 0.5 -   SrAuH3 *           |
     +---------------------------------------> omega_log (K)
         0    400   800   1200   1600
            (heavy-X bottleneck)  (light-X / light-cation only)
\end{verbatim}
Every campaign data point falls in one of two regions --- $\lambda > 2$ but
\emph{unstable}, or stable but $\omega_{\log} < \SI{700}{\kelvin}$. The
region above the $T_c = \SI{200}{\kelvin}$ isocontour is the binary
family's blind spot.

\paragraph{Margin 1: anharmonic $\lambda$-suppression ratio $S$.}
SSCHA leaves $\eta$ and $M$ fixed (the electronic structure is unchanged)
and shifts only $\langle\omega^2\rangle \to \langle\omega^2\rangle_{\mathrm{harm}}
+ \Delta\omega^2$. The numerator cancels in the ratio, giving a closed form:
\begin{equation}
S = \frac{\lambda_{\mathrm{anharm}}}{\lambda_{\mathrm{harm}}}
  = \frac{\langle\omega^2\rangle_{\mathrm{harm}}}
         {\langle\omega^2\rangle_{\mathrm{harm}} + \Delta\omega^2},
\qquad 0 < S \le 1.
\end{equation}
For the \ce{H3O} anchor ($\langle\omega^2\rangle_{\mathrm{harm}} = 1$,
$\Delta\omega^2 = 1.479$), $S = 0.4034$, so the harmonic $\lambda = 2.479$
renormalises to $\approx 1.0$ --- the centre of the observed
0.52--1.48 anharmonic band. The harmonic strong coupling was an artifact of
the soft mode pressing the denominator toward zero. Verbatim (g5):
\begin{lstlisting}[basicstyle=\ttfamily\footnotesize,frame=single,xleftmargin=0pt]
$ hexa verify --expr lambda_anharm_suppress 1.0 1.479 0.4033884626
  calc   = 0.403388  ~= expected 0.403388  (|D|=4.89956e-10 <= eps=1e-9)
  tier   = GREEN SUPPORTED-NUMERICAL  (hexa-native libm-class recompute)
\end{lstlisting}

\paragraph{Margin 2: stability--coupling margin $m$.}
The stiffness needed to support a target coupling is
$\langle\omega^2\rangle_\lambda = \eta/(M\lambda_{\mathrm{target}})$. The
signed slack of the stabilised lattice over that requirement is
\begin{equation}
m = \frac{\langle\omega^2\rangle_{\mathrm{anharm}}
         - \langle\omega^2\rangle_\lambda}
        {\langle\omega^2\rangle_{\mathrm{anharm}}},
\qquad
\begin{cases}
m > 0: & \text{escape (stable \emph{and} strong)}\\
m < 0: & \text{trapped (binary hydride)}\\
m = 0: & \text{exactly on the wall.}
\end{cases}
\end{equation}
\ce{CaH6} escapes ($m = +0.5$); \ce{H3O} is trapped ($m = -1.479$).
Verbatim (g5):
\begin{lstlisting}[basicstyle=\ttfamily\footnotesize,frame=single,xleftmargin=0pt]
$ hexa verify --expr stability_coupling_margin 1.0 0.5 0.5
  calc = 0.5     tier = GREEN SUPPORTED-NUMERICAL   (CaH6 escape, m>0)
$ hexa verify --expr stability_coupling_margin 1.0 2.479 -1.479
  calc = -1.479  tier = GREEN SUPPORTED-NUMERICAL   (H3O trapped, m<0)
\end{lstlisting}

\paragraph{Reading.}
The N5 wall is the reason the binary route ceilings near \SI{158}{\kelvin}
(stable). The escape direction is to decouple $\eta$ from
$\langle\omega^2\rangle$ --- a cation that stiffens the lattice (raising
$\langle\omega^2\rangle$ for stability) while preserving the H-derived
$N(E_F)$ (keeping $\eta$ high for coupling). \ce{CaH6}'s $m = +0.5$ is the
existence proof; whether ambient $X_2MH_6$ instances realise it is the
question Appendix~\ref{app:E} answers (negatively, so far). The two margins
above are NOVEL primitives promoted to the \texttt{hexa-lang} stdlib; the
mechanism claim that ``a cation preserves $\eta$ while raising
$\langle\omega^2\rangle$'' is honestly fenced as \tierWhite\
SPECULATION-FENCED (a DFT/wet-lab question), while only the closed-form
margins are \tierGreen.
